\section*{第 \ref{ch:g8:winds-global-circulation} 章——风与大气环流}

\begin{solution}{exo:g8:winds-global-circulation:1}
(a) 东北\omterm{def:g8:winds-global-circulation:global-winds}{信风}，来自东北。(b) \omterm{def:g8:winds-global-circulation:global-winds}{西风带}，
来自西北。(c) \omterm{def:g8:winds-global-circulation:global-winds}{极地东风带}，来自东方（从
东北方向吹来）。(d) 东南\omterm{def:g8:winds-global-circulation:global-winds}{信风}，来自东南。
\end{solution}

\begin{solution}{exo:g8:winds-global-circulation:2}
(a) 暖空气上升，冷空气下沉。(b) 低压。(c) 从高压
流向低压。(d) 向右。
\end{solution}

\begin{solution}{exo:g8:winds-global-circulation:3}
(a) 错：\omterm{def:g8:winds-global-circulation:breezes}{海风}从\omterm{def:g2:water-on-earth:ocean}{海洋}吹向陆地（白天）。
(b) 对。(c) 错：它们从东方吹来（从东北或
东南方向吹来）。(d) 对。
\end{solution}

\begin{solution}{exo:g8:winds-global-circulation:4}
图中画出一个环流：沿地面，箭头由\omterm{def:g2:water-on-earth:ocean}{海洋}指向
陆地；陆地上方，箭头向上；高空中，箭头由陆地
飘向\omterm{def:g2:water-on-earth:ocean}{海洋}；\omterm{def:g2:water-on-earth:ocean}{海洋}上方，箭头向下。H 标在
下沉气流（\omterm{def:g2:water-on-earth:ocean}{海洋}上空）之下的地面，L 标在
上升气流（陆地上空）之下的地面。
\end{solution}

\begin{solution}{exo:g8:winds-global-circulation:5}
(a) 面积是原来的 1.41 倍。(b) 50~\%。(c) 26~\%。(d) 靠近
极地时，阳光斜射并散布在较大面积上，因此
每平方米获得的能量较少。
\end{solution}

\begin{solution}{exo:g8:winds-global-circulation:6}
(a) \omterm{def:g8:winds-global-circulation:global-winds}{极地东风带}（北纬 \ang{70} 至 \ang{60}），\omterm{def:g8:winds-global-circulation:global-winds}{西风带}（北纬 \ang{60} 至
\ang{30}），东北\omterm{def:g8:winds-global-circulation:global-winds}{信风}（北纬 \ang{30} 至\omterm{def:g3:continents-and-oceans:equator}{赤道}），
东南\omterm{def:g8:winds-global-circulation:global-winds}{信风}（\omterm{def:g3:continents-and-oceans:equator}{赤道}至南纬 \ang{30}），\omterm{def:g8:winds-global-circulation:global-winds}{西风带}（南纬 \ang{30} 至
\ang{60}），\omterm{def:g8:winds-global-circulation:global-winds}{极地东风带}（南纬 \ang{60} 至 \ang{70}）。(b) 靠近
\omterm{def:g3:continents-and-oceans:equator}{赤道}，即\omterm{def:g3:continents-and-oceans:equator}{赤道}无风带。(c) 靠近北纬 \ang{30} 和南纬 \ang{30}，位于
下沉气流之下。
\end{solution}

\begin{solution}{exo:g8:winds-global-circulation:7}
(a) 陆地（$31 > 22$）：风由\omterm{def:g2:water-on-earth:ocean}{海洋}吹向陆地，是
\omterm{def:g8:winds-global-circulation:breezes}{海风}。(b) \omterm{def:g2:water-on-earth:ocean}{海洋}（$21 > 17$）：风由陆地吹向
\omterm{def:g2:water-on-earth:ocean}{海洋}，是\omterm{def:g8:winds-global-circulation:breezes}{陆风}。(c) 陆地：$31 - 17 = 14$ 度；\omterm{def:g2:water-on-earth:ocean}{海洋}：$22 -
21 = 1$ 度。(d) 太阳光只加热薄薄一层地面，
地面升温和降温都很快；海水则把热量分布到厚厚的
一层水中，因此昼夜温度几乎不变。
\end{solution}

\begin{solution}{exo:g8:winds-global-circulation:8}
(a) 从北方吹来。(b) 向右偏。(c) 从东北方向吹来（即
东北\omterm{def:g8:winds-global-circulation:global-winds}{信风}）。(d) 若无自转，风将从南方
吹来；自转使其向左偏转，因此风从东南方向吹来。
\end{solution}

\begin{solution}{exo:g8:winds-global-circulation:9}
\omterm{def:g3:continents-and-oceans:equator}{赤道}附近，太阳使地面和\omterm{def:g2:water-on-earth:ocean}{海洋}受到强烈的\emph{加热}；
湿热的空气\emph{上升}；上升时它\emph{降温}，
水汽凝结成\emph{云}，形成阵雨和
雷暴。在北纬 \ang{25} 附近，\omterm{def:g3:continents-and-oceans:equator}{赤道}上升的空气
\emph{下沉}；其下方地面气压\emph{高}；
下沉的空气增温，云消散，该地区保持\emph{干燥}：
形成\omterm{def:g5:weather-and-climate:desert}{沙漠}。
\end{solution}

\begin{solution}{exo:g8:winds-global-circulation:10}
(a) \omterm{def:g8:winds-global-circulation:global-winds}{西风带}。不对：它们从西方吹来，与向西航行的船
相向。(b) 要到达\omterm{def:g8:winds-global-circulation:global-winds}{信风}带，需位于约北纬 \ang{30}
以南。(c) 东北\omterm{def:g8:winds-global-circulation:global-winds}{信风}，位于\omterm{def:g3:continents-and-oceans:equator}{赤道}与约北纬 \ang{30}
之间。(d) \omterm{def:g8:winds-global-circulation:global-winds}{西风带}，位于约北纬 \ang{30} 以北（靠近北纬 \ang{40}，
即母港所在\omterm{def:g6:latitude-longitude:latitude}{纬度}）。
\end{solution}

\begin{solution}{exo:g8:winds-global-circulation:11}
(a) $1 \div 0.5 = 2$：表中给出 2。(b) $1 \div 0.866 \approx
1.1547$，即 \textbf{1.15}。(c) 在 \ang{30}：$1000 \div 1.15 \approx
\qty{870}{W}$，每平方米（精确值 $1000 \times 0.866 =
\qty{866}{W}$）。在 \ang{60}：$1000 \div 2 = \qty{500}{W}$。
\end{solution}

\begin{solution}{exo:g8:winds-global-circulation:12}
(a) 三个环流圈：空气在\omterm{def:g3:continents-and-oceans:equator}{赤道}上升，在 \ang{30} 附近下沉；
在 \ang{60} 附近上升，在 \ang{30} 附近下沉；在 \ang{60} 附近上升，在
极点下沉。(b) L 位于\omterm{def:g3:continents-and-oceans:equator}{赤道}和 \ang{60} 附近；H 位于 \ang{30} 附近和
极点。(c) 从 \ang{30} 到\omterm{def:g3:continents-and-oceans:equator}{赤道}，地表空气向南流动，
并向右偏转：形成东北\omterm{def:g8:winds-global-circulation:global-winds}{信风}。从 \ang{30} 到
\ang{60}，空气向北流动并向右偏转：风从
西南方向吹来。从极点到 \ang{60}，空气向南流动，向右偏转
：形成极地东风。(d) 在中间环流圈中，地表空气流向
极地，即向北；向右偏转后，风向
东北方向。这就是\omterm{def:g8:winds-global-circulation:global-winds}{西风带}。
\end{solution}

\begin{solution}{pb:g8:winds-global-circulation:1}
\textbf{1.} \omterm{def:g6:latitude-longitude:longitude}{经度}差 $59.6 - 23.5 = 36.1$ 度。

\textbf{2.} $36.1 \times 108 \approx 3899$，约 \qty{3900}{km}。

\textbf{3.} 是的：$3900 - 3896 = 4$~km，差值很小。

\textbf{4.} 是的：两地都位于北纬 \ang{14} 附近，介于\omterm{def:g3:continents-and-oceans:equator}{赤道}与约
北纬 \ang{30} 之间，即东北\omterm{def:g8:winds-global-circulation:global-winds}{信风}带。

\textbf{5.} 从东北方向吹来（箭头指向
西南）。

\textbf{6.} 向前，并略有侧向：风从船尾
稍偏右舷，把船推向西方。

\textbf{7.} 这是大气环流的一部分：空气持续在
北纬 \ang{30} 附近下沉，在\omterm{def:g3:continents-and-oceans:equator}{赤道}附近上升，因此地表空气持续
流向\omterm{def:g3:continents-and-oceans:equator}{赤道}，地球自转总是使其向
右偏。

\textbf{8.} \omterm{def:g8:winds-global-circulation:global-winds}{信风}会迎面吹向船只。船长
应向北航行，越过约北纬 \ang{30}，寻找\omterm{def:g8:winds-global-circulation:global-winds}{西风带}。

\textbf{9.} 每天 $7 \times 24 = \qty{168}{km}$。

\textbf{10.} $3 \times 168 = \qty{504}{km}$。

\textbf{11.} $3896 \div 168 \approx 23.2$ 天。

\textbf{12.} 约 \textbf{23 天}：不到一个月。
\end{solution}
